\section{Introduction}

A real Hadamard matrix of order \(n\) is a matrix \(H\) with entries
in \(\{-1,1\}\) such that \(HH^{\mathsf T}=nI_n\). It is
\emph{circulant} if there are signs \(h_0,\ldots,h_{n-1}\) for which
\[
 H_{ij}=h_{j-i\bmod n}\qquad(0\leq i,j<n).
\]
Thus every row is a cyclic shift of the first. For such a sign matrix,
define the periodic autocorrelations by
\[
 P_h(t)=\sum_{j=0}^{n-1}h_jh_{j+t\bmod n}
 \qquad(0\leq t<n).
\]
Here \(P_h(0)=n\), and the Hadamard condition is equivalent to
\(P_h(t)=0\) for every \(1\leq t<n\).
The circulant Hadamard conjecture, traditionally attributed to Ryser
\cite{Ryser1963,LeungSchmidt2012}, asserts that the only possible
orders are \(1\) and \(4\). We prove this conjecture.

\begin{theorem}\label{thm:main}
For a positive integer \(n\), a real circulant Hadamard matrix of
order \(n\) exists if and only if \(n\in\{1,4\}\).
\end{theorem}

The theorem also completes the classification of Barker-sequence
lengths. For an integer \(n>1\), a sign sequence
\(a=(a_0,\ldots,a_{n-1})\in\{-1,1\}^n\) is a \emph{Barker sequence}
if its nontrivial aperiodic autocorrelations satisfy
\[
 |C_a(t)|\leq1\quad(1\leq t<n),\qquad
 C_a(t)=\sum_{j=0}^{n-t-1}a_ja_{j+t}.
\]
The Barker-sequence conjecture asserts that no such sequence has
length greater than \(13\). Turyn and Storer established the odd-length
nonexistence theorem \cite{TurynStorer1961}; we use the later proof of
Schmidt and Willms \cite{SchmidtWillms2016} for the exact list.
For even length greater
than two, the Barker inequalities imply that all nontrivial periodic
autocorrelations vanish \cite{TurynStorer1961,Turyn1965}.
Theorem~\ref{thm:main} therefore settles the remaining even case.

\begin{corollary}[Barker-sequence lengths]\label{cor:barker}
For an integer \(n>1\), a Barker sequence of length \(n\) exists if
and only if
\[
 n\in\{2,3,4,5,7,11,13\}.
\]
\end{corollary}

Section~\ref{sec:barker} gives the short periodic-autocorrelation
argument and records examples at every listed length.

\subsection*{Context and prior work}

The sign and cyclic conditions connect the problem to cyclic
difference sets. Let \(D=\{j:h_j=-1\}\subseteq\Z/n\Z\), and put
\(k=|D|\). If \(N_D(t)\) counts ordered pairs \((a,b)\in D^2\) with
\(a-b=t\), then
\[
 P_h(t)=n-4k+4N_D(t).
\]
A cyclic \((v,k,\lambda)\) difference set is a \(k\)-element subset
of a cyclic group of order \(v\) for which every nonzero difference
occurs exactly \(\lambda\) times. For a circulant Hadamard row, the identity
\((\sum_jh_j)^2=\sum_{t=0}^{n-1}P_h(t)=n\) shows that the row sum
is \(\pm2u\) when \(n=4u^2\). Complementing \(D\) when necessary
makes that sum \(2u\), and hence
\(k=2u^2-u\). The autocorrelation identity then says precisely that
\(D\) is a cyclic difference set with parameters
\[
 (v,k,\lambda)=(4u^2,\,2u^2-u,\,u^2-u).
\]
Conversely, such a difference set supplies a circulant Hadamard row
by assigning a negative sign to its elements. This equivalence holds
for every positive \(u\); the restriction to odd \(u\) is an additional
arithmetic conclusion.

Turyn studied these cyclic difference sets through the arithmetic of
cyclotomic character sums and the Fourier relations imposed by their
common coefficient array \cite[pp.~319--323]{Turyn1965}. His
nonexistence results show that an order greater than four must be
\(4u^2\), where \(u\) is odd and not a prime power
\cite[Corollary~2 to Theorem~6, p.~333, and Theorem~8 with its
Corollary, pp.~335--336]{Turyn1965}. The proof in
Section~\ref{sec:reduction} gives a streamlined group-ring form of
his descent for the \(2\)-primary coefficients.

Later work sharpened the arithmetic restrictions. For a positive
integer \(M\), let \(\zeta_M\) be a primitive \(M\)-th root of unity.
For \(X\in\Z[\zeta_M]\) with \(X\overline X=m\in\Z_{>0}\),
Bernhard Schmidt's field descent places a multiple of \(X\) by a root
of unity in a cyclotomic subfield determined explicitly by the ambient
conductor \(M\) and \(m\) \cite{Schmidt1999}. Leung and Schmidt
developed group-ring decompositions that refine these restrictions
\cite{LeungSchmidt2005}, followed by further exclusions in their
2012 work and their anti-field-descent method
\cite{LeungSchmidt2012,LeungSchmidt2016}. Logan and Mossinghoff's
2017 computation, using searches for double Wieferich prime pairs,
excluded all but \(4\,489\) integers \(n\) with
\(4<n\leq4\cdot10^{30}\) as possible circulant Hadamard orders
\cite{LoganMossinghoff2017}. These \(4\,489\) integers remained
as candidates under those tests.

Complementary work examines the coefficients more directly. Euler,
Gallardo and Rahavandrainy obtained combinatorial restrictions on
the blocks of signs in the first row
\cite{EulerGallardoRahavandrainy2016}. Approximate orthogonality has
a different range of possibilities: Steinerberger used flat Littlewood
polynomials, which have sign coefficients and uniformly controlled
modulus on the unit circle, to show that in every order there is a
real sign circulant whose singular values lie between two fixed positive
multiples of \(\sqrt n\) \cite{Steinerberger2024}. The conjecture
requires exact equality of all singular values to \(\sqrt n\).
Complete proofs of the conjecture have also been claimed in
\cite{OhHashi2016,OrozcoLopez2019,Morris2023,Gallardo2024,ManjhiKumar2025}.

The comparison of character values also has precedents in Turyn's
work. His preliminary fact (4) obtains a root-of-unity quotient from
equality of principal ideals and absolute values. His Theorem~5 fixes
a character and compares the values obtained by multiplying it by
characters in a subgroup. It assumes equality of their principal ideals
and that the subgroup order is coprime to both the difference-set norm parameter \(k-\lambda\)
and the order of the fixed character
\cite[p.~321, item~(4), and Theorem~5, pp.~328--329]{Turyn1965}.

\subsection*{Group-ring notation}

For a finite abelian group \(V\) and a subring \(R\subseteq\C\), the
group ring \(R[V]\) consists of sums \(f=\sum_{z\in V}f_z z\), with
multiplication induced by the group law. Scalars denote multiples of
the identity element. When \(R\) is stable under complex conjugation,
put
\[
 f^*=\sum_{z\in V}\overline{f_z}\,z^{-1}.
\]
We write \(C_N\) for the cyclic group of order \(N\), and use
\(R[C_N]=R[X]/(X^N-1)\) when a generator \(X\) is chosen.
The \emph{augmentation} \(f(1)\) sends every group element to \(1\).
A character evaluation sends group elements to roots of unity, and a
quotient of groups induces a projection of group rings. Partial
character evaluation commutes with the involution, which conjugates
the evaluated coefficients on the remaining group ring.

All coefficient rings below are subrings of \(\C\), hence domains.
A localization \(R_{\m}\) at a maximal ideal allows denominators in
\(R\setminus\m\), and its residue field is
\(R_{\m}/\m R_{\m}\). The notation \(\overline z\) always denotes
complex conjugation; residues will be expressed by congruences modulo
the indicated maximal ideal.

\subsection*{Proof strategy}

Represent the first row by \(h=\sum_jh_jX^j\in\Z[C_n]\).
Orthogonality gives the scalar norm identity \(hh^*=n\).
Section~\ref{sec:reduction} first proves the classical order
restriction \(n=4u^2\) with \(u\) odd. Its cyclotomic basis turns
divisibility of a primitive character value into divisibility of
the coefficients of its remainder; at the prime two, repeated
projection then rules out a larger power of two in the order.

Suppose \(u>1\), put \(P=C_{u^2}\), and let \(I\) be the set of primes
dividing \(u\). For each \(p\in I\), choose a generator of the
\(p\)-primary component and a primitive \(p\)-th root \(\rho_p\).
Use the character sending that generator to \(\rho_p\), and put
\(B=\Z[i,\{\rho_p:p\in I\}]\). For \(S\subseteq I\), let \(x_S\)
denote the value of \(x\in\Z[i][P]\) at those characters for
\(p\in S\) and at the trivial character on the remaining components.
The norm \(xx^*=u^2\) gives \(|x_S|=u\), so the alternating product
\[
 \Delta(x)=\prod_{S\subseteq I}x_S^{(-1)^{|S|}}\in\Frac(B)
\]
is defined.

Orthogonality of translates over any finite abelian group likewise gives
a scalar group-ring norm. Here each \(p\in I\)
divides the scalar \(u^2\), and the \(p\)-component of \(P\) is cyclic
of order \(p^e\) with \(e=v_p(u^2)\). Lemma~\ref{lem:comparison}
uses this match. At each descent step, projection divides the group
order by \(p\); division of one projected factor by \(p\) then lowers
the scalar valuation by one. Both character ratios are preserved.
It follows that
\(x_{S\cup\{p\}}/x_S\) is a local unit of residue one at every maximal
ideal of \(B\) containing \(p\). Pairing the factors of \(\Delta(x)\)
in each prime direction controls it at those ideals; the norm identity
controls the remaining ideals. Thus \(\Delta(x)\) is an algebraic
integer. The alternating exponents sum to zero, so its images under
every unital homomorphism \(B\to\C\) have absolute value one.
Kronecker's criterion makes it a root of unity, and the local residues
restrict its order to a power of every \(p\in I\). In particular, that
order is odd. Proposition~\ref{prop:alternating} applies to every
\(x\in\Z[i][P]\) satisfying \(xx^*=u^2\).

The cyclic factor \(C_4\) supplies the remaining link with the signs. Write
\(C_{4u^2}=C_4\times P\), with \(Z\) a generator of \(C_4\).
Define the three normalized partial evaluations by
\[
 c=\tfrac12h\vert_{Z=1},\qquad
 d=\tfrac12h\vert_{Z=-1},\qquad
 g=\tfrac12h\vert_{Z=i}.
\]
The common sign coefficients make them integral in \(\Z[i][P]\),
and each has norm \(u^2\). Thus
\(\Delta(d)/\Delta(c)\) and \(\Delta(g)/\Delta(c)\) have odd order.

Section~\ref{sec:binary} uses the same coefficients at a maximal
ideal above two, where the odd norm makes every \(c_S\) a unit.
The binary coefficients place the ratios \(d_S/c_S\) and \(g_S/c_S\)
in the forms \(1+2L_{r,S}\) and \(1+(1+i)L_{w,S}\), with both
\(L\)-terms integral in one local ring. Hence the two alternating
ratio products have residue one; their odd order makes them exactly
one. The first-order terms of exact products have alternating sum
zero in the common residue field. But the coefficient identity gives
\[
 L_{r,S}+L_{w,S}\equiv J_S/c_S,
 \qquad J=\sum_{z\in P}z.
\]
Every nontrivial character annihilates \(J\), while its trivial value
is the odd unit \(u^2\). The resulting alternating sum therefore has
one nonzero term, which is the contradiction.
